% SLIDE 12: Training at Scale
\begin{frame}{Training at Scale}
\centering
\begin{tikzpicture}
  % Stats boxes
  \node[pipeblock=garnet, minimum width=3.2cm, minimum height=1.2cm] (imgs) at (0,0)
    {\textbf{16M}\\images};
  \node[pipeblock=atlantic, minimum width=3.2cm, minimum height=1.2cm] (ds) at (4.5,0)
    {\textbf{16}\\datasets};
  \node[pipeblock=congaree, minimum width=3.2cm, minimum height=1.2cm] (cam) at (9,0)
    {\textbf{10K+}\\camera models};

  % GPU
  \node[pipeblock=rose, minimum width=3.2cm, minimum height=1.2cm] (gpu) at (4.5,-2)
    {\textbf{48 $\times$ A100}\\GPUs};

  % Annotation
  \node[font=\small\sffamily, color=bk70, below=0.4cm of gpu, align=center] {
    Indoor: NYUv2, ScanNet, Taskonomy, \ldots\\
    Outdoor: KITTI, DDAD, Waymo, Lyft, \ldots\\
    Mixed: DIML, 3D Ken Burns, BlendedMVS, \ldots
  };
\end{tikzpicture}

\vspace{0.5cm}
\textbf{Critical data gap:} Only $\sim$20K outdoor images have GT surface normals (mostly from ScanNet indoor).
This motivates the depth-normal consistency loss $L_{d\text{-}n}$ as a way to learn normals from depth-only outdoor data.
\end{frame}
